\section{Discussion}
\label{sec:discussion}

\paragraph{Longer contexts lower the cliff.} Equation~\eqref{eq:cliff} with $\kappa$ between 1.2 and 2.2 gives a rough count
of the agents a card can hold before its cache stops paying. For a 32B model in BF16 and sessions of 68k tokens, the
median reported for GitHub Copilot's agent traces~\cite{copilot2026wild}, it gives about 4--8 sessions on an H200, 7--13 on a
192\,GB card and 13--24 on a 288\,GB card. These are projections, not measurements: our measurements support the
model at about 11k and 27k tokens on one card. If agent contexts keep growing faster than HBM capacity, the number of agents a card can hold will fall, and host tiers and
admission control will matter more.

\paragraph{Energy, capacity and money.} For an operator who rents GPUs, the energy saved is worth little in itself:
at its 600\,W maximum board power~\cite{nvidia2024h200} an H200 NVL draws about 440\,kWh a month, about \$44 at an assumed \$0.10 per kWh, while the card we rented
cost \$3.71 an hour, about \$2,700 a month. What the results buy is capacity. Finishing every production request
instead of 85--86\%, or serving 22--33\% more recorded agent turns, amounts to needing fewer GPUs. Where power, not
GPUs, limits a deployment, lower energy per turn means more turns from the same power.

\paragraph{Limitations.}
\label{sec:limits}
All results come from one GPU type (H200 NVL) in one rented container (the 10-minute adaptive-gate runs in a second container), with one model family. Energy is GPU board
energy only: host CPU and memory were not metered, and system-level savings can be proportionally smaller, since on one agent server GPU telemetry
missed 41--45\% of the energy~\cite{luo2026where}. The pauses of the recorded trajectories are synthetic, because the dataset has
no tool timings; the production trace has real timing but synthetic token contents, which preserve sharing and length
but not, for example, any effect of content on decoding. One GPU cannot carry the production trace, so we replayed
1/8 and 1/4 samples by prefix group. The admission results use
synthetic sessions or recorded trajectories, not a live agent framework, and the adaptive gate needs to know how many
sessions are open; its switch was tested only with 5\,s pauses. The miss-penalty sweep, BFCL, most production cells, the long-context no-cache baseline, the 10-minute adaptive-gate runs and SGLang's GPU cache
with Qwen3-8B ran one round. In the production replay, energy per prompt token counts all energy but only finished
requests' tokens, which favors policies that finish more requests; the 15--17\% host-tier gain includes this effect,
which we did not separate. One block of runs was
discarded after two copies of a run script shared the GPU; no result in this paper comes from it.

\section{Conclusion}

Below the cliff the GPU cache saved 38--64\% of GPU energy per turn on synthetic sessions with 5\,s pauses; past it, energy per turn was within
0.01--1.6\% of running without a cache, because a thrashing cache does the same work as no cache. The cliff follows from capacity and session size, with a measured factor of 1.2--2.2 that was larger with 30\,s than with 5\,s pauses, and
moves inward as contexts grow. A 1/4 sample of production tool-agent traffic pushed one GPU past its capacity, and a 160\,GB host tier
let it finish every request: 16--18\% more prompt tokens for about the same GPU board energy. Restoring KV took
1/110 to 1/220 of the above-idle GPU energy of recomputing it, holding it cost no measurable power, and a gate that
acts only under pressure cost no energy below the cliff and, in two 10-minute runs above it, served every session with
42--50\% less energy per turn and 48--62\% more turns, at the price of minutes of queueing at admission. On a busy GPU that keeps up, the energy effect of a host tier
can be estimated from hashed request logs before a server is changed: within 4--6\% for the two tiers we tested on production traffic.
